\documentclass[10pt]{beamer}
\usepackage{tikz}
\usepackage{xcolor}

\setbeamercolor{normal text}{fg=black}
\setbeamercolor{frametitle}{fg=white}
\setbeamercolor{title}{fg=black}

\title{Background Colors}
\subtitle{BeamerForge Global Reference}

\begin{document}

% ============================================================
\section{Light Backgrounds}

\begin{frame}{Off-White (#F5F5F5)}
\usebackgroundtemplate{
  \begin{tikzpicture}[remember picture, overlay]
    \fill[HTML]{F5F5F5} (current page.north west) rectangle (current page.south east);
  \end{tikzpicture}
}
\begin{block}{Recommended}
  The preferred background for most presentations.\\
  Soft on the eyes, works in bright rooms.\\
  Good contrast with dark text.
\end{block}
\end{frame}

\begin{frame}{Light Gray (#EEEEEE)}
\usebackgroundtemplate{
  \begin{tikzpicture}[remember picture, overlay]
    \fill[HTML]{EEEEEE} (current page.north west) rectangle (current page.south east);
  \end{tikzpicture}
}
\begin{block}{Alternative}
  Slightly darker than off-white.\\
  Softer appearance.\\
  Good for longer presentations.
\end{block}
\end{frame}

\begin{frame}{Pure White (#FFFFFF)}
\usebackgroundtemplate{
  \begin{tikzpicture}[remember picture, overlay]
    \fill[HTML]{FFFFFF} (current page.north west) rectangle (current page.south east);
  \end{tikzpicture}
}
\begin{block}{Maximum Contrast}
  Highest contrast with dark text.\\
  Use in controlled lighting.\\
  Can be harsh in bright rooms.
\end{block}
\end{frame}

% ============================================================
\section{Summary}

\begin{frame}{Quick Reference}
\centering
\begin{tabular}{l l}
\toprule
\textbf{Name} & \textbf{Hex} \\
\midrule
Off-White & \#F5F5F5 (Recommended) \\
Light Gray & \#EEEEEE (Alternative) \\
Pure White & \#FFFFFF (Maximum) \\
\bottomrule
\end{tabular}
\end{frame}

\end{document}